\documentclass[11pt]{article}
\title{Matching a shared frame is not enough when the available writer has the wrong axes}
\date{26 November 2026}
\begin{document}
\maketitle
\paragraph{Question.}
A robot owner may request a horizontal world vector $h$ and a
building operator a vertical world vector $v$. They may agree on
the same numeric frame declaration, which might not match the
pose consumed by physics. We test this shared-setup error and a
closely related control limitation: even with the correct frame
observed, does each existing command footprint represent the
effects that actor is authorized to supply?
These are declared owner obligations, not an access-control API
implemented by Gazebo. The independent contributions are modeled
at the assembly boundary before sending one physical command.

\paragraph{Exact representation condition.}
Let $H$ be the horizontal plane, $V$ its vertical orthogonal
complement and $R$ the consumer's proper rotation. Under a
declared native split the robot writer supplies local $H$ and
the lift writer supplies local $V$. When the selected phase order
preserves both, the two representable scoped world spaces are
$RH$ and $RV$. Thus a given desired pair is separately
representable exactly when $R^{-1}h\in H$ and $R^{-1}v\in V$.
All such pairs are representable exactly when $RV=V$, equivalently
$RH=H$. Even in that case the old numeric local coordinates
cannot be carried forward unchanged unless the required inverse
images stay unchanged. A positive nonzero vertical effect is not
representable by the existing lift-only slot if $RV$ is not the
world vertical line. No choice of a better scalar for the same
slot resolves that limitation. The stock velocity producer can
write all axes; restricting its input to the local plane is
part of this experiment's separate-authority policy, not a
Gazebo feature or claim about all alternative controllers.

\paragraph{Paired source controls.}
The source fixture starts with the selected native velocity
writer, lift writer, and world setter from the prior replay.
The authored identity pose and no-collision fixed geometry make
the local split work. Reuse the same shared identity assumption
at the same pose position, but give the stand-in consumer a
quarter-turn about the horizontal $y$ axis. The same requested
horizontal $x$ goes to negative world $z$ and the lift's local
vertical command goes to world $x$. The native-shaped outcome
matches the desired sum there only if desired $x$ and $z$ are both
zero (the $y$ direction is fixed).

The separate positive control grants one composed writer authority
to submit the checked sum of both contributions. It reads the
consumer-bound frame, rotates each independently authorized world
contribution back into local coordinates, verifies both components,
then writes their sum. The same selected native consumer gives
the correct total without broadening either owner's declared
world remit. The comparison exercises 54 input and frame cases.
An authored version guard leaves the initially empty slot untouched
if its observed frame is not the frame connected to consumption.
This requires the test's serial binding; it is not a solution to
infer pose or to retire an old preexisting effect.

\paragraph{Boundary.}
This supplies a scoped shared-setup and inter-owner frame-interface
comparison using selected native code and authored poses and rights.
It does not run Gazebo dynamics, test a real frame deployment,
or claim a robot or lift can operate in a tilted pose. A version
guard could be less conservative if it used the exact demand
signatures in the adjacent quotient result. A matched direct
assembler with the same wider write permission and current frame
also succeeds. No field behavior, safety certification, operating
deadline, or architecture disposition is asserted.
\end{document}
